\documentclass[11pt,a4paper]{article}
\usepackage[T1]{fontenc}
\usepackage[utf8]{inputenc}
\usepackage{amsmath,amsthm,mathtools}
\usepackage{newpxtext,newpxmath}
\usepackage[margin=25mm,headheight=15pt]{geometry}
\usepackage{microtype}
\usepackage{xcolor,booktabs,array,longtable}
\usepackage[most]{tcolorbox}
\usepackage{enumitem}
\usepackage{fancyhdr}
\usepackage{titlesec}
\usepackage{listings}
\usepackage{hyperref}
\usepackage[nameinlink,noabbrev]{cleveref}
\definecolor{ink}{HTML}{19354A}
\definecolor{accent}{HTML}{226E78}
\definecolor{pale}{HTML}{EEF5F6}
\hypersetup{colorlinks=true,linkcolor=ink,citecolor=accent,urlcolor=accent,
 pdftitle={A proof of Barry's Hankel--Somos-6 conjecture},
 pdfauthor={AI-assisted research note},pdfsubject={Hankel determinants, generating functions, Somos recurrences}}
\titleformat{\section}{\Large\bfseries\color{ink}}{\thesection}{0.65em}{}
\titleformat{\subsection}{\large\bfseries\color{ink}}{\thesubsection}{0.65em}{}
\pagestyle{fancy}\fancyhf{}
\fancyhead[L]{\small\color{ink}Barry's Hankel--Somos-6 conjecture}
\fancyhead[R]{\small Research note}
\fancyfoot[C]{\small\thepage}
\setlength{\parskip}{4pt}
\setlength{\parindent}{1.15em}
\setlength{\emergencystretch}{1.5em}
\numberwithin{equation}{section}
\newtheorem{theorem}{Theorem}[section]
\newtheorem{lemma}[theorem]{Lemma}
\newtheorem{proposition}[theorem]{Proposition}
\newtheorem{corollary}[theorem]{Corollary}
\theoremstyle{definition}
\newtheorem{definition}[theorem]{Definition}
\newtheorem{remark}[theorem]{Remark}
\newcommand{\Z}{\mathbb Z}
\newcommand{\Q}{\mathbb Q}
\newcommand{\K}{\mathbb K}
\newcommand{\coef}[2]{[#1^{#2}]}
\newcommand{\Norm}{\operatorname{N}}
\newcommand{\ord}{\operatorname{ord}}
\newcommand{\Hank}{\operatorname{H}}
\newcommand{\ct}{\operatorname{ct}}
\newcommand{\Res}{\operatorname{Res}}
\lstset{basicstyle=\small\ttfamily,breaklines=true,columns=fullflexible,
 frame=single,rulecolor=\color{accent},backgroundcolor=\color{pale},
 showstringspaces=false,xleftmargin=0.3em,xrightmargin=0.3em}
\setcounter{tocdepth}{1}

\begin{document}
\thispagestyle{empty}
\begin{center}
{\small\sffamily\color{accent}GENERATING FUNCTIONS\quad /\quad INTEGER SEQUENCES\quad /\quad EXACT PROOFS}\par
\vspace{1.1cm}
{\LARGE\bfseries\color{ink}A proof of Barry's\par Hankel--Somos-6 conjecture}\par
\vspace{0.45cm}
{\large Polynomial specialization, a homogeneous family,\par and reproducible exact verification}\par
\vspace{0.7cm}
{AI-assisted research note}\par
{19 September 2026}
\end{center}
\vspace{0.7cm}
\begin{abstract}
Conjecture 7 of Paul Barry's 2022 paper proposes a three-parameter family of
quadratic algebraic generating functions whose Hankel determinants satisfy a
Somos-6 recurrence with its middle coefficient zero. We prove the claimed
bilinear identity and its homogeneous four-parameter form. The decisive
observation is a constant square completion: the apparent cubic remainder in
the associated sextic discriminant becomes linear. This identifies the
recurrence coefficients with the invariants in van der Poorten's 2005
continued-fraction theorem. We give an algebraic derivation of the needed
local identity and the Hankel correspondence.

A genuine technical issue is that the original family has an exceptional
initial complete quotient, and some parameter values produce zero Hankel
determinants. We handle both by embedding the family in a polynomial
seven-parameter deformation. An explicit Catalan specialization proves that
all its generic Hankel determinants are nonzero. The resulting identity is
then specialized without division. Consequently the final statement holds
over every commutative coefficient ring, including at singular parameters.
The accompanying programs independently compute determinants and verify
10,248 recurrence residuals in 1,128 test cases, with additional example and
symbolic checks.
\end{abstract}
\vspace{0.4cm}
\begin{tcolorbox}[colback=pale,colframe=accent,boxrule=0.6pt,arc=1mm]
\textbf{Main coefficient simplification.} For Barry's parameters,
\[
 \alpha=t^2(r+2)^2,\qquad
 \gamma=t^3\bigl((s+t)(r+s+t+2)^2+t(r+2)^3\bigr).
\]
The second expression expands to the polynomial in Conjecture 7.
\end{tcolorbox}
\vfill
\begingroup\noindent\small\textbf{Status and attribution.}
This is an unrefereed research manuscript, not a machine-checked formalization.
The genus-two continued-fraction identity is due to van der Poorten; it is not
claimed as a new result here. The contribution is the explicit reduction,
the complete specialization argument, and a proof of the named conjecture.
A targeted literature search did not locate an earlier direct resolution of
Conjecture 7, but does not establish a priority claim.\par\endgroup
\newpage
\tableofcontents
\vspace{0.6cm}
\begin{tcolorbox}[colback=white,colframe=accent,boxrule=0.4pt,title=Reading guide,
 coltitle=white,colbacktitle=accent]
The problem and answer are in Sections 1--2. Section 4 isolates the algebraic
observation that makes the problem tractable. Sections 5--7 contain the
all-index proof, including the handling of zero determinants. Sections 8--10
provide examples, consequences, and a reproducibility audit.
\end{tcolorbox}

\section{The selected problem and its status}

Let $G(x)=\sum_{n\geq0}g_nx^n$ be the formal power series with $G(0)=1$ determined by
\begin{equation}\label{eq:barry}
 G(x)=\frac{1}{1-x\dfrac{1+rx}{1-x}
                    -x^2\dfrac{1+sx}{1-x}-tx^3G(x)}.
\end{equation}
Barry \cite[Conjecture 7, p.~7]{Barry2022} conjectured that the Hankel transform
of its coefficients is an $(\alpha,0,\gamma)$ Somos-6 sequence. The absence of
the middle term means that the identity has only three terms, despite its
index width being six. In the general Somos-6 recurrence there is an
additional term $\beta H_{N-2}H_{N-4}$; the conjecture asserts $\beta=0$.

We use matrix size as the determinant index:
\begin{equation}\label{eq:hdef}
 H_0=1,\qquad H_N=\det(g_{i+j})_{0\leq i,j<N}\quad(N\geq1).
\end{equation}
Barry's indexing is $h_n=H_{n+1}$. Keeping this shift explicit avoids a common
source of mistakes in comparisons of Hankel recurrences.

\begin{theorem}[Barry's Conjecture 7, in polynomial form]\label{thm:barry}
For the series \eqref{eq:barry}, the determinants \eqref{eq:hdef} satisfy
\begin{equation}\label{eq:barryrec}
 H_NH_{N-6}=\alpha H_{N-1}H_{N-5}+\gamma H_{N-3}^2
 \qquad(N\geq6),
\end{equation}
where
\begin{equation}\label{eq:barrycoeff}
 \alpha=t^2(r+2)^2,\qquad
 \gamma=t^3\bigl((s+t)(r+s+t+2)^2+t(r+2)^3\bigr).
\end{equation}
This is an identity in $\Z[r,s,t]$ and hence remains true under every
specialization into a commutative ring. Over a field, whenever $H_{N-6}\ne0$,
\eqref{eq:barryrec} is the usual quotient-form Somos recurrence.
\end{theorem}

With $h_n=H_{n+1}$, the conjectured identity follows for $n\geq6$. The theorem
also supplies the preceding boundary relation corresponding to $N=6$, with
the empty determinant $H_0=1$.

\subsection{What was checked in the literature}
The source problem is a specific displayed conjecture, not an artificially
constructed question. Wang and Zhang \cite{WangZhang2024} proved the four
Somos-4 conjectures numbered 2--5 in Barry's paper; their Corollaries 4--7
identify those numbers explicitly. That result should not be mistaken for a
proof of the Somos-6 conjecture considered here.

The central earlier theorem is van der Poorten's width-six relation
\cite[Theorems 3.1--3.2]{vdP2005}. Hone's treatment
\cite[Section 4 and Theorem 5.4]{Hone2021} places continued fractions and
Hankel determinants in a broader hyperelliptic framework. Those results
made the problem a plausible target: the missing task was to identify the
particular curve and determinant normalization, rather than to create an
entire theory of Somos sequences.

The search was checked against recent related work as well. The stated
applications in Chang--Liu \cite{ChangLiu2026} concern bilateral Somos-4 and
Somos-5, and Scheuerle's September 2026 preprint \cite{Scheuerle2026} concerns
Somos-4. No direct resolution of Barry's Conjecture 7 was located. This is a
bounded literature-search statement, not a certification that no unpublished
or differently phrased proof exists. The proof below stands independently
of that bibliographic uncertainty.

\section{A homogeneous formulation of the answer}

It is convenient to replace $r+2$ and $s+t$ by independent parameters and to
allow a scale parameter $u$. Define
\begin{equation}\label{eq:D}
 D(x)=1-2ux+(u^2-R)x^2+(t-B)x^3
\end{equation}
and let $F(x)=\sum_{n\geq0}f_nx^n$, $f_0=1$, solve
\begin{equation}\label{eq:mainGF}
 D(x)F(x)=1-ux+tx^3(1-ux)F(x)^2.
\end{equation}
The coefficient of the new unknown $f_n$ is always $1$, so this equation
uniquely defines $F$ over $\Z[u,R,B,t]$, or over any ring after specialization.
Write $H_N(F)=\det(f_{i+j})_{0\leq i,j<N}$ and $H_0(F)=1$.

\begin{theorem}[Homogeneous Hankel--Somos-6 identity]\label{thm:main}
For all $N\geq6$,
\begin{equation}\label{eq:mainrec}
 \begin{split}
 H_N(F)H_{N-6}(F)
 ={}&t^2R^2H_{N-1}(F)H_{N-5}(F)\\
 &+t^3\bigl(B(B+uR)^2+tR^3\bigr)H_{N-3}(F)^2.
 \end{split}
\end{equation}
The equality holds in $\Z[u,R,B,t]$ without any nonvanishing hypotheses.
\end{theorem}

The substitution
\begin{equation}\label{eq:originalspec}
 u=1,\qquad R=r+2,\qquad B=s+t
\end{equation}
turns \eqref{eq:mainGF} into \eqref{eq:barry}, after multiplication by $1-x$.
Thus \cref{thm:main} proves \cref{thm:barry}.

Some useful initial values are
\begin{align}
 f_0&=1,\qquad f_1=u,\qquad f_2=u^2+R,\notag\\
 f_3&=u^3+3uR+B,\qquad f_4=u^4+6u^2R+R^2+3uB.
\end{align}
Direct determinants give
\begin{align}\label{eq:initialH}
 H_0&=H_1=1,\qquad H_2=R,\qquad H_3=-B(B+uR),\\
 H_4&=-t\bigl(B(B+uR)^2+tR^3\bigr).\notag
\end{align}
These formulas already show why it would be incorrect to assume that all
specialized determinants are nonzero or positive.

The fourth parameter is a homogeneous extension, not a claim of an
independent new moduli direction: when $u$ is invertible, it can be removed
by scaling $x$. Its inclusion makes the weights and the exceptional case
$u=0$ transparent.

\section{Formal coefficients and a combinatorial interpretation}

\subsection{A triangular coefficient recurrence}
Set $f_j=0$ for $j<0$ and
\[
 C_m(F)=\sum_{j=0}^{m}f_jf_{m-j}\quad(m\geq0),\qquad C_m(F)=0\quad(m<0).
\]
Equating coefficients in \eqref{eq:mainGF} yields
\begin{equation}\label{eq:momrec}
 \begin{split}
 f_n={}&\mathbf1_{n=0}-u\mathbf1_{n=1}+2uf_{n-1}
 +(R-u^2)f_{n-2}+(B-t)f_{n-3}\\
 &\hspace{2em}+tC_{n-3}(F)-utC_{n-4}(F).
 \end{split}
\end{equation}
For $n\geq1$, only previously determined coefficients occur. In particular,
all moments and all Hankel determinants are integer polynomials in the
parameters. This elementary observation will ultimately permit division-free
specialization of the more complicated continued-fraction argument.

\subsection{A finite Catalan formula}
Let $C(z)=\sum_{j\geq0}C_jz^j$ be the Catalan series, characterized formally by
$C(z)=1+zC(z)^2$, where $C_j=\frac{1}{j+1}\binom{2j}{j}$. Substitution into
\eqref{eq:mainGF} gives
\begin{equation}\label{eq:catalanF}
 F(x)=\frac{1-ux}{D(x)}\,
 C\!\left(\frac{tx^3(1-ux)^2}{D(x)^2}\right).
\end{equation}
No square-root branch or analytic convergence convention is needed. The
Catalan form in Barry's original problem is also displayed in
\cite[Conjecture 7]{Barry2022}.

For example, \eqref{eq:catalanF} gives the finite coefficient formula
\begin{equation}\label{eq:finitecoef}
 f_n=\sum_{j=0}^{\lfloor n/3\rfloor}C_jt^j
       [x^{n-3j}]\frac{(1-ux)^{2j+1}}{D(x)^{2j+1}}.
\end{equation}
One can remove coefficient extraction entirely. In the following expression,
$m=2j+1$ and $(m)_L=m(m+1)\cdots(m+L-1)$, with $(m)_0=1$:
\begin{equation}\label{eq:multinomial}
 \begin{split}
 f_n={}&\sum_{j=0}^{\lfloor n/3\rfloor}C_jt^j
 \sum_{\substack{0\leq q\leq m,\ \ell_1,\ell_2,\ell_3\geq0\\
 q+\ell_1+2\ell_2+3\ell_3=n-3j}}
 \binom{m}{q}(-u)^q
 \frac{(m)_{\ell_1+\ell_2+\ell_3}}{\ell_1!\ell_2!\ell_3!}\\
 &\hspace{9em}\cdot(2u)^{\ell_1}(R-u^2)^{\ell_2}(B-t)^{\ell_3}.
 \end{split}
\end{equation}
This follows by expanding the numerator binomially and the denominator by
the negative multinomial theorem. The displayed factorial ratio is an
integer multinomial coefficient. Thus this is also a direct finite formula
for the integer sequences in the family.

\subsection{Colored unary--binary trees for the original parameters}
The original equation can be written in a particularly simple form:
\begin{equation}\label{eq:treegrammar}
 G=1+W(x)G+tx^3G^2,\qquad
 W(x)=x+(r+2)x^2+(r+s+2)\frac{x^3}{1-x}.
\end{equation}
For nonnegative integers $r,s,t$, interpret a tree recursively as an empty
object, a unary root with one subtree, or a binary root with two ordered
subtrees. A binary root has size $3$ and one of $t$ colors. A unary root of
size $1$ has one color; one of size $2$ has $r+2$ colors; one of any size
$\ell\geq3$ has $r+s+2$ colors. The total size is the sum of root sizes.
All nonempty roots have positive size, so every coefficient counts a finite
set of objects. The recursive construction has exactly
\eqref{eq:treegrammar} as its generating equation.

This gives a concrete combinatorial source of the moments. It does not make
the Hankel determinants positive: determinants involve signs, and the
negative values in \eqref{eq:initialH} are entirely compatible with positive
moment counts. The verification code also expands \eqref{eq:treegrammar}
independently and compares it with \eqref{eq:momrec}.

\section{The decisive square completion}

Put $X=x^{-1}$ and introduce the cubic
\begin{equation}\label{eq:S}
 S(X)=X^3-2uX^2+(u^2-R)X-B+t.
\end{equation}
The discriminant of the quadratic equation for $F$ corresponds to the
sextic
\begin{equation}\label{eq:sextic}
 Y^2=S(X)^2-4tX(X-u)^2.
\end{equation}
At first sight, the remainder next to the square is cubic. The appropriate
square, however, is not $S^2$. Set
\begin{equation}\label{eq:A}
 A(X)=S(X)-2t
 =X^3-2uX^2+(u^2-R)X-B-t.
\end{equation}
A direct expansion gives
\begin{equation}\label{eq:squarecomplete}
 \boxed{\ S(X)^2-4tX(X-u)^2=A(X)^2-4t(RX+B).\ }
\end{equation}
Indeed, $S(X)-X(X-u)^2=-RX-B+t$, so subtracting the two
squares leaves $-4t(RX+B)$.

Generically, put $v=tR$ and $w=-B/R$. Then the same curve is
\begin{equation}\label{eq:curve}
 Y^2=A(X)^2-4v(X-w),\qquad
 Z=\frac{A+Y}{2},\qquad Z^2-AZ+v(X-w)=0.
\end{equation}
The invariant coefficients in the width-six continued-fraction identity are
$v^2$ and $-v^3A(w)$. In our case,
\begin{equation}\label{eq:invariants}
 v^2=t^2R^2,\qquad
 -v^3A(w)=t^3\bigl(B(B+uR)^2+tR^3\bigr).
\end{equation}
This explains both the compact coefficient formula and the disappearance of
the middle Somos-6 term. The curve is generically of genus two, but no
smoothness assumption is required in the final polynomial identity.

There is also an explicit identification of the moment function. Choose
$Y=A+O(X^{-2})$ at infinity (the square-root branch with leading term $X^3$), and put $M(X)=X^{-1}F(X^{-1})$. The quadratic
formula, or direct substitution, gives
\begin{equation}\label{eq:Mcurve}
 M(X)=\frac{S-Y}{2t(X-u)}=\frac{X(X-u)}{Z+t}.
\end{equation}
For the second equality, use the norm identity
\begin{equation}\label{eq:initialnorm}
 (Z+t)(\overline Z+t)=t(A+t)+t(RX+B)=tX(X-u)^2.
\end{equation}
Here the bar denotes the other root of the quadratic equation; it is not
complex conjugation.

The natural initial complete quotient is therefore
\[
 \frac1M=\frac{Z+t}{X(X-u)}.
\]
Its numerator correction is constant. In the customary notation
$P_0=d_0(X+e_0)$, this is a limiting situation with $d_0=0$ rather than an
ordinary generic initial value. One must not simply invoke a theorem that
assumes all the $d_j$ are nonzero and then silently ignore this issue.
Moreover, $w=-B/R$ is undefined when $R=0$. The next two sections resolve
both problems by proving a polynomial identity in a larger family first.

\section{A polynomial Hankel form of the genus-two identity}

Work initially over the polynomial ring
\[
 \mathcal S=\Z[a,b,c,p,k,\lambda,\mu].
\]
Define
\begin{equation}\label{eq:univpolys}
 A=X^3+aX^2+bX+c,\qquad P=pX+k,\qquad Q=X^2+\lambda X+\mu.
\end{equation}
Divide the polynomial $P(A+P)$ by the monic polynomial $Q$ and write the
result as
\begin{equation}\label{eq:division}
 P(A+P)+vX+m=QL.
\end{equation}
Thus $-(vX+m)$ is the remainder, and $L=pX^2+\ell_1X+\ell_0$ is the quotient.
Everything in \eqref{eq:division} is polynomial over $\mathcal S$. For
implementation, convenient formulas are
\begin{align}\label{eq:univcoeffs}
 \ell_1&=ap+k-\lambda p,\\
 \ell_0&=ak-a\lambda p+bp-k\lambda+\lambda^2p-\mu p+p^2,\notag\\
 v&=\ell_1\mu+\ell_0\lambda-cp-bk-2pk,\notag\\
 m&=\ell_0\mu-ck-k^2.\notag
\end{align}

Let $\Phi(x)=\sum_{n\geq0}\phi_nx^n$, $\phi_0=1$, be defined by
\begin{equation}\label{eq:univGF}
 \begin{split}
 \bigl(1+ax+(b+2p)x^2+(c+2k)x^3\bigr)\Phi
  ={}&1+\lambda x+\mu x^2\\
  &+x^2(p+\ell_1x+\ell_0x^2)\Phi^2.
 \end{split}
\end{equation}
Equivalently, $\mathcal M(X)=X^{-1}\Phi(X^{-1})$ is the unique Laurent series
$X^{-1}+O(X^{-2})$ satisfying
\begin{equation}\label{eq:univM}
 Q-(A+2P)\mathcal M+L\mathcal M^2=0.
\end{equation}

\begin{theorem}[Polynomial Hankel--Somos identity]\label{thm:universal}
Let $K_N=\det(\phi_{i+j})_{0\leq i,j<N}$ and $K_0=1$. For every $N\geq6$,
\begin{equation}\label{eq:univH}
 K_NK_{N-6}=v^2K_{N-1}K_{N-5}+\Gamma K_{N-3}^2,
\end{equation}
where
\begin{equation}\label{eq:Gamma}
 \Gamma=m^3-avm^2+bv^2m-cv^3.
\end{equation}
The identity is valid in $\mathcal S$, and hence after any commutative-ring
specialization. No nonvanishing assumption belongs to the conclusion.
\end{theorem}

The generic continued-fraction identity behind this theorem is van der
Poorten's \cite{vdP2005}. The polynomial family and the argument below make
its use at exceptional initial data explicit. In particular, the proof does
not assume, without verification, that Barry's smaller family meets a
suitable generic open set.

\section{Proof of the polynomial theorem}

\subsection{A single specialization proves generic nonvanishing}
Every $\phi_n$, and therefore every $K_N$, belongs to $\mathcal S$ by the
triangular coefficient recursion in \eqref{eq:univGF}. To prove that none of
these polynomials is identically zero, specialize
\begin{equation}\label{eq:catalanspec}
 (a,b,c,p,k,\lambda,\mu)=(0,-3,0,1,0,0,-1).
\end{equation}
Here $L=X^2-1$, $v=0$, and $m=1$. Equation \eqref{eq:univGF} becomes
\[
 (1-x^2)\Phi=(1-x^2)+x^2(1-x^2)\Phi^2,
\]
so $\Phi=C(x^2)$.

For completeness, all Hankel determinants of the interleaved Catalan
sequence are $1$. Let $J$ be the adjacency matrix of the path on vertices
$0,1,2,\ldots$. Products of this infinite matrix are harmless here: all
vectors used have finite support. A walk of length $2j$ from $0$ back to $0$
is a Dyck path, so
\[
 (J^{2j})_{00}=C_j,\qquad (J^{2j+1})_{00}=0.
\]
Define monic polynomials $\pi_0=1$, $\pi_1=X$, and
$\pi_{n+1}=X\pi_n-\pi_{n-1}$. Induction gives $\pi_n(J)e_0=e_n$.
For the moment functional $\mathcal L(q)=e_0^{\mathsf T}q(J)e_0$, symmetry of
$J$ now gives
\[
 \mathcal L(\pi_i\pi_j)=\delta_{ij}.
\]
The change from monomials to the monic $\pi_i$ is unit lower triangular.
Consequently the Gram determinant in the monomial basis is also $1$.
This proves $K_N=1$ under \eqref{eq:catalanspec}, for every $N$.

It follows that each $K_N$ is nonzero in the fraction field
$\mathcal K=\Q(a,b,c,p,k,\lambda,\mu)$. Notice that the specialization has
$v=0$; that is no obstacle. Its sole purpose is to prove that the Hankel
polynomials are nonzero. The polynomial $v$ itself is not identically zero,
as is evident from \eqref{eq:univcoeffs}, so $v$ is invertible in $\mathcal K$.
No positivity assertion is being made for generic moments.

\subsection{The Jacobi coefficients are Hankel ratios}
We recall the precise formal correspondence, with the indexing required here.
Suppose $\mathcal M=X^{-1}\sum_{j\geq0}\phi_jX^{-j}$ and all $K_N$ are nonzero.
Define a functional by $\mathcal L(X^j)=\phi_j$. The linear equations for the
monic degree-$n$ polynomial orthogonal to $1,X,\ldots,X^{n-1}$ have coefficient
matrix with determinant $K_n$, so they have a unique solution $\pi_n$.
If $\nu_n=\mathcal L(\pi_n^2)$, a monic change of basis in the Gram determinant
gives
\[
 \nu_n=\frac{K_{n+1}}{K_n}.
\]
Multiplication by $X$ is symmetric for this bilinear form. Orthogonality and
degree therefore imply a three-term recurrence
\begin{equation}\label{eq:orthrec}
 \pi_{n+1}=(X-\beta_n)\pi_n-b_n\pi_{n-1},\qquad
 b_n=\frac{\nu_n}{\nu_{n-1}}
     =\frac{K_{n-1}K_{n+1}}{K_n^2}\quad(n\geq1).
\end{equation}
There is no analytic measure assumption.

The same recurrence gives the convergents to the formal Jacobi fraction
\begin{equation}\label{eq:Jfrac}
 \mathcal M=
 \cfrac{1}{X-\beta_0-
   \cfrac{b_1}{X-\beta_1-
   \cfrac{b_2}{X-\beta_2-\cdots}}}.
\end{equation}
For detail, let $\rho_n$ be the polynomial part of $\mathcal M\pi_n$.
Then $\rho_0=0$ and $\rho_1=1$. For $n\geq1$, the coefficient of $X^{-1}$
in $\mathcal M\pi_n$ is $\mathcal L(\pi_n)=0$. Taking polynomial parts in
\eqref{eq:orthrec} therefore shows that the $\rho_n$ satisfy the same
three-term recurrence as the $\pi_n$. Consequently $\rho_n/\pi_n$ are the
continued-fraction convergents in \eqref{eq:Jfrac}.
Furthermore, all coefficients of $X^{-1},\ldots,X^{-n}$ in
$\mathcal M\pi_n$ vanish by orthogonality. Dividing the remaining tail by
the monic polynomial $\pi_n$ shows that
$\mathcal M-\rho_n/\pi_n=O(X^{-2n-1})$.
Agreement to arbitrarily high order proves the formal identity.
This also establishes that all polynomial partial quotients are linear.
These are the classical formulas discussed in \cite[Theorem 4.1]{Hone2021};
the argument above fixes the index shift explicitly.

\subsection{Polynomial complete quotients}
Set
\begin{equation}\label{eq:Zuniv}
 Z=\frac{Q}{\mathcal M}-P.
\end{equation}
Using \eqref{eq:division} and \eqref{eq:univM}, direct substitution gives
\begin{equation}\label{eq:Zquad}
 Z^2-AZ+vX+m=0.
\end{equation}
The chosen branch has $Z=A+O(X^{-2})$; its conjugate is
$\overline Z=A-Z=O(X^{-2})$. The initial quotient is
\[
 F_0=\frac1{\mathcal M}=\frac{Z+P_0}{Q_0},\qquad P_0=P,\quad Q_0=Q.
\]
Ordinary polynomial continued-fraction division produces
\begin{equation}\label{eq:CFstep}
 F_n=\frac{Z+P_n}{Q_n}=a_n+\frac1{F_{n+1}},\qquad
 P_{n+1}=a_nQ_n-A-P_n,
\end{equation}
with
\begin{equation}\label{eq:CFnorm}
 Q_nQ_{n+1}=-\bigl(P_{n+1}(A+P_{n+1})+vX+m\bigr).
\end{equation}
The first norm divisibility is \eqref{eq:division}. Subsequent divisibility
follows inductively by reducing the norm modulo $Q_n$ and using
$A+P_{n+1}\equiv-P_n\pmod{Q_n}$.

Starting with $Q_0$ of degree two and $P_0$ of degree one, the formula for
$P_{n+1}$ cancels the terms of degrees three and two, so $\deg P_{n+1}\leq1$.
The norm formula then gives $\deg Q_{n+1}\leq2$. Since the next quotient has
numerator of degree three and has a linear polynomial part, its denominator
must have degree exactly two. This proves the degree assertions inductively.
Write
\begin{equation}\label{eq:normal}
 P_n=d_n(X+e_n),\qquad
 Q_n=u_nq_n,\qquad q_n=X^2+\lambda_nX+\mu_n,
 \qquad u_0=1.
\end{equation}
Here the $u_n$ are local normalization constants and are unrelated to the
homogeneous parameter $u$ in \cref{thm:main}. Comparison of leading
coefficients in \eqref{eq:CFnorm} gives
\begin{equation}\label{eq:unprod}
 u_{n-1}u_n=-d_n\qquad(n\geq1).
\end{equation}
In particular $d_n\ne0$. If a numerator correction had vanished in degree
one, the next denominator would have degree below two and the next partial
quotient would not be linear. This cannot occur in the generic moment
problem just established.

The polynomial part is
$a_n=(X+a-\lambda_n)/u_n$. Rescaling adjacent levels of the continued
fraction with \eqref{eq:unprod} changes its numerators to $-d_n$. Comparison
with \eqref{eq:Jfrac} therefore yields the essential identification
\begin{equation}\label{eq:dH}
 \boxed{\quad d_n=\frac{K_{n-1}K_{n+1}}{K_n^2}\quad(n\geq1).\quad}
\end{equation}
For example, the beginning of the rescaled fraction is
\[
 \mathcal M=\cfrac{1}{X+a-\lambda_0-
            \cfrac{d_1}{X+a-\lambda_1-
            \cfrac{d_2}{X+a-\lambda_2-\cdots}}}.
\]
There is no extra factor of $2$ or $4$ in this normalization, because $Z$,
not $Y=2Z-A$, is used throughout.

\subsection{A local proof of van der Poorten's identity}
Set $w=-m/v$, so $vX+m=v(X-w)$. We now derive the needed width-six identity
using only polynomial division and evaluations. This is the identity in
\cite[Theorem 3.1]{vdP2005}, with the $X^2$ coefficient of $A$ retained.

The normalized Euclidean and norm relations are
\begin{align}\label{eq:E}
 P_n+P_{n+1}&=(X+a-\lambda_n)q_n-A,\\
 d_nq_{n-1}q_n&=P_n(A+P_n)+v(X-w).\label{eq:N}
\end{align}
Reducing the second relation modulo $q_n$, and using the first, shows that
$P_nP_{n+1}-v(X-w)$ is divisible by $q_n$. Its degree is two and its leading
coefficient is $d_nd_{n+1}$. Hence
\begin{equation}\label{eq:cross}
 d_nd_{n+1}q_n=P_nP_{n+1}-v(X-w).
\end{equation}
This simple remainder identity is the main algebraic tool.

We first show that
\begin{equation}\label{eq:qeval}
 q_n(-e_n)=-d_{n-1}.
\end{equation}
For full detail, the coefficients of $X^3$ and $X^2$ in \eqref{eq:N} give
\begin{align}
 \lambda_{n-1}+\lambda_n&=a+e_n,\label{eq:coeff3}\\
 \lambda_{n-1}\lambda_n+\mu_{n-1}+\mu_n&=b+ae_n+d_n.\label{eq:coeff2}
\end{align}
The coefficient of $X$ in \eqref{eq:E} at index $n-1$ is
\begin{equation}\label{eq:coeff1}
 d_{n-1}+d_n=\mu_{n-1}+a\lambda_{n-1}-\lambda_{n-1}^2-b.
\end{equation}
Eliminating $\mu_{n-1}$ between these equations gives
\[
 d_{n-1}=(a-\lambda_{n-1})e_n-\mu_n
         =\lambda_ne_n-e_n^2-\mu_n=-q_n(-e_n),
\]
as claimed.

Evaluate \eqref{eq:cross} at $X=-e_n$. Since $P_n(-e_n)=0$,
\eqref{eq:qeval} yields
\begin{equation}\label{eq:triple}
 T_n:=d_{n-1}d_nd_{n+1}=-v(e_n+w).
\end{equation}
All the $d_j$ in use are nonzero, so $e_n+w\ne0$ in the fraction field.
On the other hand, evaluating \eqref{eq:cross} at $X=w$ gives
\begin{equation}\label{eq:qw}
 q_n(w)=(e_n+w)(e_{n+1}+w).
\end{equation}
The norm relation \eqref{eq:N}, also evaluated at $w$ and divided by $d_n$,
is
\[
 q_{n-1}(w)q_n(w)=(e_n+w)\bigl(A(w)+d_n(e_n+w)\bigr).
\]
Substituting \eqref{eq:qw} and canceling the nonzero factor $e_n+w$ gives
\begin{equation}\label{eq:ecubic}
 (e_{n-1}+w)(e_n+w)(e_{n+1}+w)=A(w)+d_n(e_n+w).
\end{equation}
Multiply by $-v^3$ and use \eqref{eq:triple} at three consecutive indices.
The result is
\begin{equation}\label{eq:Tlocal}
 T_{n-1}T_nT_{n+1}=v^2d_nT_n-v^3A(w),
\end{equation}
or, explicitly,
\begin{equation}\label{eq:vdp}
 d_{n-2}d_{n-1}^2d_n^3d_{n+1}^2d_{n+2}
 =v^2d_{n-1}d_n^2d_{n+1}-v^3A(w).
\end{equation}
This proof explains why retaining the $X^2$ term in $A$ causes no problem:
only coefficient comparisons and evaluations of the full cubic have been
used. For our subsequent application, $n\geq3$, so no negative-index
continued-fraction data are required.

\subsection{Telescoping and removal of all denominators}
Substitution of \eqref{eq:dH} gives two telescoping identities:
\begin{align}\label{eq:telescope1}
 d_{n-1}d_n^2d_{n+1}&=\frac{K_{n-2}K_{n+2}}{K_n^2},\\
 d_{n-2}d_{n-1}^2d_n^3d_{n+1}^2d_{n+2}
 &=\frac{K_{n-3}K_{n+3}}{K_n^2}.\label{eq:telescope2}
\end{align}
Multiplying \eqref{eq:vdp} by $K_n^2$ and setting $N=n+3$ gives the desired
Hankel relation in $\mathcal K$, with coefficient $-v^3A(w)$. But
\begin{equation}\label{eq:clearinvariant}
 -v^3A(-m/v)=m^3-avm^2+bv^2m-cv^3=\Gamma
\end{equation}
is already polynomial. Thus both sides of \eqref{eq:univH} belong to
$\mathcal S$ and are equal in its fraction field. The injection
$\mathcal S\hookrightarrow\mathcal K$ proves their equality in
$\mathcal S$ itself.

Finally, apply any ring homomorphism from $\mathcal S$ into a commutative
ring. The specialized moments are still the unique coefficients defined by
\eqref{eq:univGF}, since its triangular recursion never divided by a
parameter. Determinants and all terms of \eqref{eq:univH} specialize as
polynomials. This proves the theorem even when $v$, $p$, a normalization
constant, or one or more Hankel determinants become zero. It also proves
validity in positive characteristic.\hfill$\square$

\section{Specialization: completion of the conjecture's proof}

In the universal theorem make the substitution
\begin{equation}\label{eq:bigSpec}
 \begin{gathered}
 a=-2u,\qquad b=u^2-R,\qquad c=-B-t,\\
 p=0,\qquad k=t,\qquad\lambda=-u,\qquad\mu=0.
 \end{gathered}
\end{equation}
Then
\[
 A=X^3-2uX^2+(u^2-R)X-B-t,\qquad P=t,\qquad Q=X(X-u).
\]
The division identity becomes
\[
 t(A+t)+tRX+tB=tX(X-u)^2,
\]
so
\begin{equation}\label{eq:Lspec}
 L=t(X-u),\qquad v=tR,\qquad m=tB.
\end{equation}
Putting these expressions into \eqref{eq:univGF} gives precisely
\eqref{eq:mainGF}.

The invariant \eqref{eq:Gamma} specializes to
\begin{align}
 \Gamma
 &=t^3\bigl(B^3+2uRB^2+(u^2-R)R^2B+(B+t)R^3\bigr)\notag\\
 &=t^3\bigl(B^3+2uRB^2+u^2R^2B+tR^3\bigr)\notag\\
 &=t^3\bigl(B(B+uR)^2+tR^3\bigr).\label{eq:gammaspec}
\end{align}
Equation \eqref{eq:univH} is therefore \eqref{eq:mainrec}. This proves
\cref{thm:main}, including the exceptional specialization $p=0$.
Using \eqref{eq:originalspec} proves \cref{thm:barry}.

To compare with the exact expression printed by Barry, expansion of the
bracket in \eqref{eq:barrycoeff} gives
\begin{equation}\label{eq:expandedgamma}
 \begin{split}
 \gamma/t^3={}&r^3t+r^2(s+7t)
 +2r\bigl(s^2+2(t+1)s+t(t+8)\bigr)\\
 &+s^3+s^2(3t+4)+s(t+2)(3t+2)+t(t^2+4t+12).
 \end{split}
\end{equation}
The notation on the left means the polynomial bracket, not a required
division when $t=0$. This is the coefficient in Conjecture 7. Thus the result
is a proof of the proposed formula, not a nearby recurrence with adjusted
parameters.

For another normalization check, the first complete quotient in the smaller
family has polynomial part $X-u$ and
\[
 P_1=RX+B,\qquad Q_1=-(RX+B)(X-u).
\]
Its leading numerator coefficient is $d_1=R$, exactly as required by
$H_0H_2/H_1^2=R$. This check is not needed for the specialization proof, but
it catches possible sign and index errors in the curve identification.

\section{Examples, singular cases, and a source correction}

\subsection{A positive moment sequence with signed Hankel determinants}
For $r=s=t=1$,
\[
 G(x)=1+x+4x^2+12x^3+34x^4+103x^5+309x^6+934x^7+\cdots.
\]
The determinants, with the empty determinant included, start as
\[
 1,\ 1,\ 3,\ -10,\ -77,\ -807,\ 437,\ 468332,\ 2665431,\ 183243907,\ldots.
\]
The coefficients are $(\alpha,\gamma)=(9,77)$. The first two instances of
\eqref{eq:barryrec} are checked by
\begin{align*}
 437&=9(-807)+77(-10)^2,\\
 468332&=9\cdot437\cdot3+77(-77)^2.
\end{align*}
These checks illustrate the normalization; the theorem, not these examples,
proves every subsequent relation.

\begin{table}[htbp]
\centering\small
\caption{Some original-family examples; entries are exact direct determinants.}
\label{tab:examples}
\begin{tabular}{@{}rrrrrrrr@{}}
\toprule
$r$&$s$&$t$&$\alpha$&$\gamma$&$H_2$&$H_3$&$H_4$\\
\midrule
0&0&1&4&17&2&$-3$&$-17$\\
1&1&1&9&77&3&$-10$&$-77$\\
$-3$&$-2$&1&1&$-5$&$-1$&$-2$&5\\
$-2$&$-2$&1&0&$-1$&0&$-1$&1\\
2&2&2&64&3072&4&$-32$&$-768$\\
\bottomrule
\end{tabular}
\end{table}

\subsection{What happens when a determinant vanishes}
At $(r,s,t)=(-2,-2,1)$, the directly computed beginning is
\[
 (H_N)_{N\geq0}=1,1,0,-1,1,0,-1,-1,0,1,-1,\ldots,
\]
and the identity is
\[
 H_NH_{N-6}=-H_{N-3}^2.
\]
For $N=8$, the expression $H_8=(-H_5^2)/H_2$ would be $0/0$ and is not a
recurrence that can be evaluated by division. The polynomial identity is
nevertheless the valid equation $0=0$. The determinants supply a canonical
sequence through this singularity; the bilinear relation alone need not
uniquely determine a continuation from singular initial data.

Likewise, when $R=0$ in the homogeneous family, $H_2=0$ and $w=-B/R$ cannot
be used, but the final coefficients $\alpha=0$ and $\gamma=t^3B^3$ are
perfectly well defined. The polynomial proof is doing real work here.

When $t=0$, $F=(1-ux)/D$ is rational. Its moments satisfy an order-at-most-three
constant-coefficient recurrence from index $3$ onward. Every fourth Hankel
column is a linear combination of the preceding three, so $H_N=0$ for
$N\geq4$. Both Somos coefficients also vanish, consistently with the theorem.

\subsection{A parameter typo in Barry's Example 8}
The source labels Example 8 with $(r,s,t)=(-2,-2,1)$ but prints a moment
sequence beginning $1,1,0,-3,-7,-9,\ldots$ and coefficients
$(\alpha,\gamma)=(1,-5)$ \cite[p.~7]{Barry2022}. Those data instead correspond
to $(r,s,t)=(-3,-2,1)$. The discrepancy is already visible in
\[
 g_2=r+3:
\]
it equals $1$ at $r=-2$ and $0$ at $r=-3$. Direct calculation at the corrected
parameters gives
\[
 (H_N)_{N\geq0}=1,1,-1,-2,5,17,-3,-122,1201,-2980,\ldots,
\]
which matches the source's displayed Hankel transform after removing $H_0$.
This is an example-label correction, not a counterexample to Conjecture 7.
The coefficient identity itself is correct.

\section{Consequences and simple transformations}

\subsection{Integrality without a Laurent-phenomenon hypothesis}
For integer parameters the moments are integers by \eqref{eq:momrec}, and
so are the Hankel determinants. Thus the quotient Somos recurrence, whenever
defined, produces integers along this determinant family. No separate
Laurent-phenomenon theorem, or assumption that all six initial values are
units, is required for this conclusion. Integrality here is immediate from
the determinant representation and should not be advertised as a new
integrality theorem for arbitrary Somos-6 initial data.

\subsection{Binomial translations preserve the entire Hankel sequence}
Let $z$ be a new parameter and define
\[
 \widetilde f_n=\sum_{j=0}^{n}\binom nj z^{n-j}f_j,
 \qquad
 \widetilde F(x)=\frac1{1-zx}F\!\left(\frac{x}{1-zx}\right).
\]
The corresponding $N\times N$ Hankel matrices are related by
\[
 \widetilde{\mathcal H}_N=T_N\mathcal H_NT_N^{\mathsf T},\qquad
 (T_N)_{ij}=\begin{cases}\binom ij z^{i-j},&j\leq i,\\0,&j>i.\end{cases}
\]
To check an entry, use the Vandermonde identity
$\sum_{p+q=j}\binom ip\binom kq=\binom{i+k}j$. Since $T_N$ has diagonal
entries $1$, all Hankel determinants are unchanged. Consequently the same
Somos recurrence holds for every binomial translate. This is a standard
Hankel invariance; the short proof is included to show exactly what happens
to this family.

\subsection{Scaling and weights}
Replacing $f_n$ by $\rho^nf_n$ multiplies $H_N$ by
$\rho^{N(N-1)}$, by factoring $\rho^i$ from row $i$ and $\rho^j$ from column
$j$. The two recurrence coefficients therefore scale by $\rho^{10}$ and
$\rho^{18}$, respectively. In \cref{thm:main}, the parameter weights are
\[
 \operatorname{wt}(u)=1,\quad \operatorname{wt}(R)=2,\quad
 \operatorname{wt}(B)=\operatorname{wt}(t)=3.
\]
Both coefficient formulas have exactly the predicted weights. This provides
an additional structural check on the exponents in the result.

More explicitly, for invertible $u$, the homogeneous series is
$F(x)=G(ux)$ with the original parameters
\[
 r=R/u^2-2,\qquad s=(B-t)/u^3,\qquad \tau=t/u^3
\]
in place of $(r,s,t)$. The polynomial theorem additionally covers $u=0$,
where this rescaling is unavailable.

\section{Reproducible exact computations}

The mathematical proof is symbolic and valid for all indices. Computation
was used to select the target, detect normalization errors, and audit the
finished formulas. It is deliberately independent of the proposed Somos
recurrence.

\subsection{What the programs actually do}
The standard-library program \texttt{code/verify.py} generates moments from
the quadratic defining equation, constructs each Hankel matrix separately,
and evaluates its determinant by fraction-free Bareiss elimination with row
pivoting. Every integer division is checked for exactness. It never generates
a test determinant from the Somos recurrence. For rational parameter tests,
exact \texttt{Fraction} arithmetic is used. Small determinant computations
are also checked against the permutation definition of the determinant.

The full deterministic run uses seed $20260919$ and the following test groups:
\begin{center}
\small\setlength{\tabcolsep}{4.5pt}
\begin{tabular}{@{}lrr@{}}
\toprule
Group & Parameter cases & Maximum matrix size\\
\midrule
Original family, $(r,s,t)\in\{-3,\ldots,3\}^3$ &343&14\\
Homogeneous family, $(u,R,B,t)\in\{-2,\ldots,2\}^4$&625&14\\
Seven-parameter polynomial family, seeded integers&128&14\\
Larger-matrix tests in the polynomial family&16&24\\
Homogeneous family with rational parameters&16&10\\
\midrule
Total&1128&24\\
\bottomrule
\end{tabular}
\end{center}

The saved run reports $10{,}248$ zero recurrence residuals in these cases.
Among the $1{,}128$ cases, $553$ contain at least one zero Hankel determinant
of size at least two; there are $3{,}465$ such zero determinant values in the
test records. Thus the test set does not silently exclude the singular
specializations that require care in the proof. There are also $84$
small-matrix determinant cross-checks, $343$ independent original-family
moment-recursion comparisons, and $25$ Catalan-specialization determinants
$H_0,\ldots,H_{24}$, all equal to $1$.

Five additional named examples are recorded through $H_{20}$ and through
moment $g_{38}$, adding $75$ recurrence checks. These are reported separately
from the $10{,}248$ main-suite residuals. The exact output and all parameters
are in \texttt{data/verification\_summary.json} and
\texttt{data/verification\_cases.json}; the example values are also supplied
as CSV files.

\subsection{Finite symbolic certificates}
The optional program \texttt{code/symbolic\_certificates.py} uses SymPy to
verify $16$ explicit algebraic identities. They include the square
completion, the initial norm factorization, the equality of the expanded
and compact gamma formulas, the universal polynomial division, all essential
specialization formulas, the coefficient elimination proving
$q_n(-e_n)=-d_{n-1}$, both Hankel telescoping identities, and the symbolic
initial determinants $H_2,H_3,H_4$. The saved run used SymPy $1.14.0$.

These finite checks are not described as a proof of an infinite recurrence.
The all-index step is the local polynomial identity in Section 6.4, followed
by telescoping and polynomial specialization. No Lean or other proof-assistant
verification is claimed.

\subsection{Commands and complexity}
From the archive root, run
\begin{lstlisting}
python code/verify.py
python code/symbolic_certificates.py
latexmk -pdf -interaction=nonstopmode -halt-on-error article.tex
\end{lstlisting}
The first command requires only Python 3.10 or later. The second requires
SymPy; an optional requirements file is included. A smaller numerical run is
available with \texttt{python code/verify.py --quick}.

The direct verification method uses $O(N^2)$ arithmetic operations to obtain
moments through order $2N$ by a quadratic convolution recursion, and
$O(N^3)$ operations for one $N\times N$ determinant. Computing every
$H_0,\ldots,H_N$ separately therefore takes $O(N^4)$ arithmetic operations.
These are operation counts, not bit-complexity estimates; exact integers can
grow rapidly. Once the theorem is proved, a nonzero Hankel sequence can be
extended much faster by its Somos recurrence, but doing that in the verifier
would make the numerical evidence circular. Singular cases should continue
to be evaluated from moments and determinants unless a separate justified
singularity-handling method is supplied.

\section{Conclusion and scope of the result}

The selected conjecture is resolved by the polynomial identity
\eqref{eq:barryrec}. The key reduction is not an accidental match of initial
terms: the sextic discriminant is exactly a cubic square plus a linear
remainder, and its two curve invariants reproduce the conjectured
coefficients for every parameter value. The seven-parameter deformation
provides a rigorous route from generic continued fractions to exceptional
initial data and zero determinants.

The result does not prove Barry's Somos-8 conjectures, classify arbitrary
Somos-6 initial-value problems, establish positivity of the Hankel
polynomials, or assert that all sequences satisfying the bilinear relation
have this determinant representation. Those would be separate questions.
The principal mathematical ingredient belongs to the existing work of
van der Poorten and the established Hankel/continued-fraction correspondence.
What this manuscript adds is an explicit application with every normalization
and specialization step made checkable.

Independent expert review remains appropriate before making a publication
or priority claim. The archive is designed to facilitate that review: it
contains the editable proof, its compiled version, a bounded literature
status note, and reproducible computations rather than only an assertion
that a conjecture has been solved.

\clearpage
\begin{thebibliography}{9}
\bibitem{Barry2022}
P. Barry, \emph{Conjectures on Somos 4, 6 and 8 sequences using Riordan arrays
and the Catalan numbers}, arXiv:2211.12637v1 (2022).
Conjecture 7 and Example 8 are on p.~7.
\url{https://arxiv.org/abs/2211.12637}.

\bibitem{vdP2005}
A. J. van der Poorten, \emph{Curves of Genus 2, Continued Fractions, and
Somos Sequences}, Journal of Integer Sequences \textbf{8} (2005), Article 05.3.4.
\url{https://cs.uwaterloo.ca/journals/JIS/VOL8/Poorten2/vdp89.pdf}.

\bibitem{Hone2021}
A. N. W. Hone, \emph{Continued fractions and Hankel determinants from
hyperelliptic curves}, Communications on Pure and Applied Mathematics
\textbf{74} (2021), 2310--2347. arXiv:1907.05204.
\url{https://arxiv.org/abs/1907.05204}.

\bibitem{WangZhang2024}
Y. Wang and Z. Zhang, \emph{Sufficient condition for $(\alpha,\beta)$ Somos 4
Hankel determinants}, Discrete Mathematics \textbf{347} (2024), 113937.
Preprint arXiv:2305.05995v2 (2023).
\url{https://arxiv.org/abs/2305.05995}.

\bibitem{ChangLiu2026}
X.-K. Chang and J. Liu, \emph{Hankel Determinants from Quadratic Orthogonal
Pairs for Hyperelliptic Functions and Their Applications},
arXiv:2603.11670v4 (2026).
\url{https://arxiv.org/abs/2603.11670}.

\bibitem{Scheuerle2026}
T. Scheuerle, \emph{Direct Generation of a Somos-4 Sequence from an Algebraic
Generating Function}, arXiv:2609.15754v1 (2026).
\url{https://arxiv.org/abs/2609.15754}.
\end{thebibliography}
\end{document}
